\begin{table*}[!t]
	\caption{Residual DC offset (\%FS of the nominal 5-V ADC span), mean $\pm$ sample standard deviation over $N{=}10$ paired baseline-seeded restarts per reported board/gain cell. Within each cell, the upper and lower values correspond to the shared fixed-reference and calibrated conditions, respectively.}
	\label{tab:results}
	\centering
	\scriptsize
	\setlength{\tabcolsep}{3.2pt}
	\renewcommand{\arraystretch}{0.82}
	
	\resizebox{\textwidth}{!}{%
		\arrayrulecolor{gray!45}
		\begin{tabular}{l cc ;{2pt/2pt} cc ;{2pt/2pt} cccc}
			\arrayrulecolor{black}
			\toprule
			Stage
			& A, $\times1$
			& A, $\times4$
			& A, $\times16$
			& B, $\times16$
			& A, $\times32$
			& B, $\times32$
			& A, $\times50$
			& B, $\times50$ \\
			\midrule
			
			PGA
			& \shortstack{$3.5\pm0.2$\\[-1.5pt]$3.4\pm0.3$}
			& \shortstack{$2.4\pm0.1$\\[-1.5pt]$2.1\pm0.3$}
			& \shortstack{$2.4\pm0.1$\\[-1.5pt]$1.9\pm0.2$}
			& \shortstack{$2.1\pm0.1$\\[-1.5pt]$1.7\pm0.1$}
			& \shortstack{$1.0\pm0.2$\\[-1.5pt]$0.5\pm0.2$}
			& \shortstack{$1.0\pm0.2$\\[-1.5pt]$0.5\pm0.2$}
			& \shortstack{$1.2\pm0.3$\\[-1.5pt]$1.0\pm0.5$}
			& \shortstack{$2.5\pm0.1$\\[-1.5pt]$1.8\pm0.3$}
			\\
			
			\arrayrulecolor{gray!45}
			\hdashline[0.4pt/1.4pt]
			\arrayrulecolor{black}
			
			BP
			& \shortstack{$1.5\pm2.2$\\[-1.5pt]$1.2\pm0.2$}
			& \shortstack{$4.9\pm0.3$\\[-1.5pt]$1.0\pm0.2$}
			& \shortstack{$3.8\pm0.9$\\[-1.5pt]$1.5\pm1.5$}
			& \shortstack{$0.66\pm0.02$\\[-1.5pt]$0.66\pm0.02$}
			& \shortstack{$7.0\pm0.4$\\[-1.5pt]$1.0\pm0.5$}
			& \shortstack{$0.66\pm0.01$\\[-1.5pt]$0.66\pm0.01$}
			& \shortstack{$1.0\pm0.7$\\[-1.5pt]$0.9\pm0.5$}
			& \shortstack{$0.65\pm0.01$\\[-1.5pt]$0.65\pm0.02$}
			\\
			
			\arrayrulecolor{gray!45}
			\hdashline[0.4pt/1.4pt]
			\arrayrulecolor{black}
			
			SUM
			& \shortstack{$\mathbf{21.8\pm1.1}^{\dagger}$\\[-1.5pt]$0.8\pm0.2$}
			& \shortstack{$6.7\pm1.0$\\[-1.5pt]$1.0\pm0.3$}
			& \shortstack{$4.9\pm1.5$\\[-1.5pt]$1.4\pm1.5$}
			& \shortstack{$4.5\pm0.8$\\[-1.5pt]$0.6\pm0.5$}
			& \shortstack{$7.4\pm0.3$\\[-1.5pt]$0.6\pm0.2$}
			& \shortstack{$8.0\pm0.8$\\[-1.5pt]$1.2\pm0.3$}
			& \shortstack{$0.4\pm0.2$\\[-1.5pt]$\mathbf{0.8\pm0.6}^{\ddagger}$}
			& \shortstack{$0.9\pm0.6$\\[-1.5pt]$0.9\pm0.6$}
			\\
			
			\arrayrulecolor{gray!45}
			\hdashline[0.4pt/1.4pt]
			\arrayrulecolor{black}
			
			LP
			& \shortstack{$\mathbf{45.8\pm3.1}^{\dagger}$\\[-1.5pt]$3.2\pm0.9$}
			& \shortstack{$\mathbf{42.6\pm4.2}^{\dagger}$\\[-1.5pt]$3.9\pm0.8$}
			& \shortstack{$\mathbf{29.6\pm6.8}^{\dagger}$\\[-1.5pt]$2.3\pm1.4$}
			& \shortstack{$3.0\pm0.1$\\[-1.5pt]$0.7\pm0.1$}
			& \shortstack{$\mathbf{12.4\pm1.3}^{\dagger}$\\[-1.5pt]$1.7\pm0.7$}
			& \shortstack{$2.9\pm0.1$\\[-1.5pt]$0.4\pm0.2$}
			& \shortstack{$4.1\pm2.1$\\[-1.5pt]$1.2\pm0.4$}
			& \shortstack{$2.9\pm0.1$\\[-1.5pt]$0.7\pm0.2$}
			\\
			
			\bottomrule
		\end{tabular}%
		\arrayrulecolor{black}
	}
	
	\vspace{1pt}
	\footnotesize
	Lower is better.
	$^{\dagger}$Fixed reference $\geq10\%$FS.
	$^{\ddagger}$The marked calibrated mean is larger than its baseline;
	equality after rounding is not marked
	(Section~\ref{sec:pga-note}).
\end{table*}

\section{Proof-of-Concept Evaluation}
\label{sec:experimental-validation}

\subsection{Initial Analog-Path Check}

With no calibrated seismic source, the Ma-derived electronics were checked independently of the geophone. An Agilent~33220A waveform generator excited the front end in overlapping batches of 0.01~Hz--200~kHz logarithmic sinusoidal sweeps, while a Tektronix TDS~2004B recorded the stage inputs and outputs. The measured PGA-output-to-ADC-input analog response was defined by the LP-output/PGA-output ratio (CH4/CH1), so it covers the electronics downstream of the monitored PGA output.

For the approximate sensor-to-ADC shape in Fig.~\ref{fig:analog_response}, the measured electronics response was cascaded with the nominal SM-24 acceleration model \cite{SM24}. All magnitude curves are normalized: the measured composite is set to 0~dB at its maximum; the ideal chain and 0--2~k$\Omega$ tuning envelope use that same reference, whereas the standalone geophone acceleration and velocity curves are each normalized to their own maximum. Local coherence, computed from the squared normalized cross-term of the analytic input and output within logarithmic sweep bins, first falls below 0.9 at approximately 0.84~kHz; the dash-dotted vertical line marks the rounded 1-kHz onset. Above this point, strong attenuation and nonmonotonic coherence prevent a reliable amplitude comparison with the model. The growing disagreement is therefore treated as measurement-limited rather than a verified change in circuit behavior. This is an initial circuit-level consistency check, neither a calibrated sensor-chain characterization nor validation of the offset-calibration method. End-to-end measurements with a calibrated seismic source remain future work.



\subsection{Prototype and Comparison Conditions}

The main comparison contains eight board--gain cells, each measured through ten paired restarts: all five PGA settings on Board~A and the $\times16$, $\times32$, and $\times50$ settings on Board~B. Thus, $N{=}10$ represents repeated restarts of the same board and gain, not ten independent devices. The shared fixed-reference baseline, $\boldsymbol{c}_{\mathrm{base}}$, was selected empirically on Board~B at $\times50$ and then used unchanged on Board~A and at the other gains. For each calibrated condition, the four stages were processed sequentially and the resulting DAC codes were frozen during acquisition. Post-calibration EEPROM writes were disabled in these measurements, so every restart began again from $\boldsymbol{c}_{\mathrm{base}}$.

Board~B at $\times1$ and $\times4$ were omitted from the eight-cell summary: their baseline signals fell within the normal deadbands, so the algorithm applied no corrective updates and no distinct after-calibration condition exists. A separate narrower-deadband diagnostic (Section~\ref{sec:pga-note}) confirmed the expected response.

\subsection{Residual DC-Offset Measurement and Results}

The LP output is the primary outcome, delivering the conditioned signal to the acquisition ADC; PGA, BP, and SUM document sequential correction and expose near-rail intermediate operating points.

At each setting, a four-channel Tektronix TDS~2004B (8-bit, 60~MHz) recorded $V_{\mathrm{mean}}$ at PGA, BP, SUM, and LP over a 1s window after the transient settled. All probe grounds were tied to $V_{\mathrm{ref}}=V_{\mathrm{DDA}}/2$, so each channel measured $V_{\mathrm{stage}}-V_{\mathrm{ref}}$; because the USB host was unearthed while the oscilloscope ground was at protective earth, this also fixed $V_{\mathrm{ref}}$ to earth potential \cite{TektronixTDS2000B}. USB-supplied $V_{\mathrm{DDA}}$ ran slightly below 5~V, and other leakage paths were not characterized. Within each restart the fixed-reference condition was recorded first and the calibrated condition second, preserving the pairing; as this order was never reversed, slow drift cannot be fully separated from the calibration effect.

Table~\ref{tab:results} reports the mean and sample standard deviation of $100|V_{\mathrm{mean}}|/5\,\mathrm{V}$ over ten paired restarts. Values near a supply rail may be clipped and indicate saturation rather than exact linear offsets. These descriptive statistics characterize within-configuration repeatability only: no oscilloscope uncertainty budget was performed, and displayed digits do not imply sub-millivolt accuracy.

Figure~\ref{fig:calibration_trace} applies the 128-tap calibration FIR to oscilloscope records resampled at the controller rate; the colored curves approximate the controller's filtered outputs but are external reconstructions, not logs, excluding native sampling, switching, resets, and fixed-point arithmetic. Table~\ref{tab:results} uses later oscilloscope records, not the controller's final ADC samples.

\begin{figure}[t]
	\centering
	\includegraphics[width=\linewidth]{imagenes/autocalibration_scope_sum.png}
	\caption{Stage deviations relative to $V_{\mathrm{ref}}$ during calibration at $\times16$ (total differential gain $\times32$). Dashed light-gray traces are oscilloscope measurements; colored curves are resampled 128-tap-FIR reconstructions, not internal logs (Section~\ref{sec:pga-note}).}
	\label{fig:calibration_trace}
\end{figure}

\begin{figure}[t]
	\centering
	\includegraphics[width=\linewidth]{imagenes/waterfall_wiggle_horizontal.png}
	\caption{Wiggle-trace field record acquired after foreground calibration; the trigger and first-arrival moveout demonstrate multichannel operation over the 10--50~m line.}
	\label{fig:waterfall}
\end{figure}

Across the eight reported board$\times$gain cells, the calibrated LP sample mean is 0.4--3.9\%FS versus 2.9--45.8\%FS under the shared fixed-reference baseline; individual calibrated observations span 0.1--5.2\%FS (Table~\ref{tab:results}). The LP mean was lower after calibration in every reported cell, without board- or gain-specific controller retuning; a one-sided sign test on the paired LP magnitudes rejects the no-reduction null in all eight cells (10 of 10 restarts lower per cell, exact $p\approx10^{-3}$).

\subsection{Field Operation Demonstration}

This field trial was an operational demonstration, not a controlled offset or surface-wave comparison. Fig.~\ref{fig:waterfall} shows a 10--50~m record with each geophone planted and kept stationary during the preceding calibration; the resolved first-arrival moveout demonstrates operation. Across laboratory and field measurements, the only condition that triggered a timeout was powering a board while its geophone was moving, whose AC transient saturated the chain and violated the quasi-static assumption.

